% !TEX program = tectonic
\documentclass[10.5pt]{article}
\usepackage{newtxtext,newtxmath}
\usepackage[a4paper,margin=0.86in]{geometry}
\usepackage{amsmath}
\usepackage{booktabs,longtable,array}
\usepackage{xcolor,enumitem,titlesec}
\usepackage[hidelinks,colorlinks=true,linkcolor=blue!55!black,urlcolor=teal!60!black]{hyperref}

\definecolor{g4}{HTML}{C65335}
\definecolor{teal}{HTML}{267B72}
\definecolor{muted}{HTML}{536A73}
\newcommand{\rv}{\mathbf r}
\newcommand{\Fv}{\mathbf F}
\newcommand{\uv}{\mathbf u}
\newcommand{\ev}{\widehat{\mathbf e}}
\newcommand{\keylabel}{\par\smallskip\noindent\textcolor{g4}{\textbf{KEY EQUATION}}\par\vspace{-5pt}}
\newcommand{\unit}[1]{\,\mathrm{#1}}
\setlist[itemize]{leftmargin=1.35em,itemsep=2pt,topsep=3pt}
\titleformat{\section}{\Large\bfseries\color{black!85}}{\thesection}{0.55em}{}
\titleformat{\subsection}{\large\bfseries\color{teal}}{\thesubsection}{0.55em}{}
\allowdisplaybreaks

\title{\vspace{-1.2cm}\textbf{The Generation-4 Single-Cell Collagen Model}\\[3pt]
\large Theory, Governing Equations, Staged Controls, and Evidence Boundaries\\[6pt]
\normalsize\textnormal{Motor--clutch collagen-remodelling project}}
\author{}
\date{2026-09-07}

\begin{document}
\maketitle
\vspace{-1cm}

\begin{abstract}
Generation 4 (G4) is a two-dimensional, single-cell mechanism-calibration platform in
micrometres, nanonewtons, and seconds. It fixes the cell first, calibrates an elastic
bead--spring collagen network, traces crosslink-mediated force transmission, adds stochastic
motor--clutch failure, and finally releases rigid-cell translation and rotation. G4 v1 is the
short-time prototype. G4 v2 extends the observation window and exposes small displacements and
fast rupture events without inflating the physical parameters. This document distinguishes
implemented equations from modelling assumptions and literature-supported mechanisms. G4 is not
yet a calibrated three-dimensional tumour-migration prediction.
\end{abstract}

\keylabel
\begin{equation}
 \zeta_b\dot{\rv}_i=
 \Fv_i^{\mathrm{stretch}}+\Fv_i^{\mathrm{bend}}+\Fv_i^{\mathrm{xlink}}
 +\Fv_i^{\mathrm{rep}}+\Fv_i^{\mathrm{active}} .
 \label{eq:backbone}
\end{equation}
Equation~\eqref{eq:backbone} is the common G4 mechanical backbone: inertia-free collagen beads
move according to elastic, contact, and active forces divided by bead drag.

\tableofcontents

\section{Positioning and stage logic}

G4 was introduced because earlier generations mixed changes in network geometry, anchoring,
contact, and cell mobility. G4 instead follows a controlled inheritance rule: once a mechanism is
accepted, the next stage changes one major ingredient while retaining the previous mechanics.

\begin{longtable}{>{\raggedright\arraybackslash}p{0.15\textwidth}p{0.29\textwidth}p{0.45\textwidth}}
\toprule
Stage & New test & Primary interpretation \\
\midrule
G4A & One-factor elastic calibration with a fixed cell & Sensitivity of displacement, traction,
radial order, and boundary response to one coefficient at a time. \\
G4B & Crosslink-graph transmission & Whether non-contact fibres move only through retained
elastic connections. \\
G4C & Independent versus shared-load clutch failure & Binding, loading, rupture cascades,
complete site loss, traction drop, and fibre recoil. \\
G4D & Release rigid-cell translation and rotation & Whether asymmetric collagen reactions are
sufficient to generate motion without prescribed polarity. \\
\bottomrule
\end{longtable}

G4 v1 runs for two to five minutes and reveals the controls, but many indirect displacements are
subpixel and many clutch events occur between saved frames. G4 v2 retains the same single-cell
concept and parameter language while using two-hour elastic and moving-cell runs, a 30-minute
clutch overview, and a 20-second event microscope sampled every 0.1 seconds.

\section{Collagen network mechanics shared by G4A--D}

Each collagen fibre is a connected chain of visible beads. Consecutive beads carry axial force;
triplets resist changes in curvature; retained intersections couple material points on different
fibres. Only beads in the outer boundary band are anchored in the main condition.

\subsection{Overdamped update}

For every non-fixed bead,
\begin{equation}
 \rv_i(t+\Delta t)=\rv_i(t)+\frac{\Delta t}{\zeta_b}\Fv_i^{\mathrm{total}},
 \qquad \rv_i(t)=\rv_i^0\quad (i\in\mathcal B_{\mathrm{fixed}}).
 \label{eq:euler}
\end{equation}
Bead drag controls the transient speed. It dissipates motion but stores no elastic energy.
The bonds and crosslinks store energy. G4 contains no standard-linear-solid element, inertia,
thermal Brownian forcing, or permanent remodelling.

\subsection{Axial stretch and compression softening}

For a bond from bead $i$ to $j$,
\begin{align}
 \mathbf d_{ij}&=\rv_j-\rv_i, & \ell_{ij}&=\lVert\mathbf d_{ij}\rVert,
 & e_{ij}&=\ell_{ij}-\ell_{ij}^0,\\
 \Fv_{i\leftarrow j}^{\mathrm{stretch}}&=k_{ij}e_{ij}\ev_{ij},
 &\ev_{ij}&=\mathbf d_{ij}/\ell_{ij},
 &k_{ij}&=
 \begin{cases}
 EA/\ell_{ij}^0,&e_{ij}\ge 0,\\
 \rho EA/\ell_{ij}^0,&e_{ij}<0.
 \end{cases}
 \label{eq:stretch}
\end{align}
The baseline compression ratio is $\rho=0.10$. Tension therefore uses the full axial stiffness,
whereas compression is softened to represent collagen microbuckling.

\subsection{Discrete bending}

The curvature stencil for an interior triplet is
\begin{equation}
 \Delta\boldsymbol\kappa_i=
 (\rv_{i-1}-2\rv_i+\rv_{i+1})-\boldsymbol\kappa_i^0,
 \qquad
 U_b=\frac12\sum_i\frac{EI}{\bar\ell_i^{,3}}
 \lVert\Delta\boldsymbol\kappa_i\rVert^2.
 \label{eq:bend}
\end{equation}
The reference curvature $\boldsymbol\kappa_i^0$ is the initial fibre shape. The bending
multiplier changes $EI$ without changing the axial rigidity $EA$.

\subsection{Permanent material-point crosslinks}

For an intersection located at fractions $\alpha_a$ and $\alpha_b$ along two host bonds,
\begin{align}
 \mathbf p_a&=(1-\alpha_a)\rv_{a_1}+\alpha_a\rv_{a_2},
 &\mathbf p_b&=(1-\alpha_b)\rv_{b_1}+\alpha_b\rv_{b_2},\\
 \Fv_x&=k_x\big[(\mathbf p_b-\mathbf p_a)-\mathbf d_x^0\big].
 \label{eq:xlink}
\end{align}
Interpolation weights distribute $\Fv_x$ back to the four host beads. The link resists
separation but is freely hinged: it does not prescribe the angle between fibres. G4 crosslinks
do not rupture or re-form.

\section{G4A: one-factor mechanics calibration}

G4A keeps the cell fixed and compares the same initial collagen geometry while varying one
coefficient at a time. The left panel is the baseline control; the right panel is the selected
variant. Both use the same cell position, candidate intersections, duration, and random seed.

Candidate intersections receive stored uniform marks $u_j$. Crosslink probability changes the
network through a nested rule:
\keylabel
\begin{equation}
 \text{retain link }j \quad\Longleftrightarrow\quad u_j<p_x,
 \qquad u_j\sim\mathcal U(0,1),
 \label{eq:nested}
\end{equation}
so increasing $p_x$ adds links without reshuffling the links already present at lower $p_x$.

\begin{longtable}{>{\raggedright\arraybackslash}p{0.24\textwidth}p{0.28\textwidth}p{0.38\textwidth}}
\toprule
Coefficient & Equation location & Expected effect when increased \\
\midrule
Total pull $F_{\mathrm{pull}}$ & Active contact force & More direct loading until stiffness or
geometry limits dominate. \\
Bending multiplier & $EI$ in Eq.~\eqref{eq:bend} & Fibres resist rotation and straightening. \\
Collagen modulus $E$ & $EA$ and $EI$ & Axial extension and usually local displacement decrease. \\
Bead drag $\zeta_b$ & Eqs.~\eqref{eq:backbone}--\eqref{eq:euler} & Transients slow; the final
elastic equilibrium should not change. \\
Crosslink probability $p_x$ & Eq.~\eqref{eq:nested} & More force paths, with possible additional
rotational constraint. \\
Crosslink stiffness $k_x$ & Eq.~\eqref{eq:xlink} & Each retained link transmits separation force
more strongly. \\
Contact width $\sigma_c$ & Eq.~\eqref{eq:weights} & The same finite pull is spread more evenly
among eligible contacts. \\
\bottomrule
\end{longtable}

The mobile-boundary control releases outer anchors only in the comparison panel. It identifies
whole-network translation that might otherwise be mistaken for local remodelling. The G4 v2
elastic overview is two hours with sparse display frames; the solver step remains
$\Delta t=0.05\unit{s}$, so display sampling does not change the physics.

\section{Hybrid cell--collagen coupling}

The closest material point on each candidate bead-to-bead segment is computed first. A point is
eligible only when it lies inside a finite surface-contact shell. Gaussian weights are normalized
only over this eligible set $\mathcal C$:
\keylabel
\begin{align}
 d_m&=\lVert\mathbf p_m-\rv_c\rVert-R,
 &0&\le d_m\le w_c,\\
 w_m&=\frac{\exp(-d_m^2/\sigma_c^2)}
 {\displaystyle\sum_{n\in\mathcal C}\exp(-d_n^2/\sigma_c^2)}.
 \label{eq:weights}
\end{align}
The baseline uses $w_c=3\unit{\mu m}$ and $\sigma_c=1.5\unit{\mu m}$. The Gaussian therefore
distributes a finite local pull; it is not evaluated over the whole ECM.

The active force and its interpolation to the two host beads are
\begin{align}
 \Fv_m^{\mathrm{active}}
 &=F_{\mathrm{pull}}w_m\frac{\rv_c-\mathbf p_m}{\lVert\rv_c-\mathbf p_m\rVert},\\
 (\Fv_i^{\mathrm{active}},\Fv_j^{\mathrm{active}})
 &=\big((1-\alpha_m)\Fv_m^{\mathrm{active}},\alpha_m\Fv_m^{\mathrm{active}}\big).
 \label{eq:active}
\end{align}
Using a material point avoids bead-resolution artefacts and preserves both total force and its
first moment along the bond.

Beads cannot penetrate the rigid circular cell:
\begin{equation}
 \delta_i=\max\big(0,R+c-\lVert\rv_i-\rv_c\rVert\big),\qquad
 \Fv_i^{\mathrm{rep}}=k_{\mathrm{rep}}\delta_i
 \frac{\rv_i-\rv_c}{\lVert\rv_i-\rv_c\rVert}.
 \label{eq:rep}
\end{equation}
This one-sided penalty is contact, not adhesion.

\section{G4B: crosslink-mediated indirect realignment}

Only direct-contact fibres receive Eq.~\eqref{eq:active}. Every non-contact fibre receives zero
Gaussian force. Fibres are classified by shortest path through the retained crosslink graph:
\keylabel
\begin{equation}
 d_{\mathrm{graph}}(f,\mathcal C_{\mathrm{direct}})\in\{0,1,2+,\infty\}.
 \label{eq:graph}
\end{equation}
Here $0$ denotes direct contact, $1$ one retained crosslink hop, $2+$ a longer connected path,
and $\infty$ an unconnected fibre. An indirect fibre moves because deformation of a direct fibre
separates the linked material points, Eq.~\eqref{eq:xlink} produces equal-and-opposite force, and
the neighbouring fibre redistributes that force through its own stretching and bending.

For nematic orientation, a fibre pointing at $\theta$ is equivalent to one pointing at
$\theta+\pi$. The implemented orientation summaries are
\begin{align}
 \theta_f&=\frac12\operatorname{atan2}
 \left(\langle\sin 2\theta_e\rangle,\langle\cos 2\theta_e\rangle\right),\\
 \Delta\theta_f&=\frac12\operatorname{atan2}
 \left(\sin 2(\theta_f-\theta_f^0),\cos 2(\theta_f-\theta_f^0)\right),\\
 u_f&=\frac{1}{N_f}\sum_{i\in f}\lVert\rv_i-\rv_i^0\rVert.
 \label{eq:metrics}
\end{align}
Each graph class reports displacement and orientation distributions. In the no-crosslink
negative control, non-contact transmitted motion should collapse toward the unconnected class.

\section{G4C: motor--clutch loading and stochastic slippage}

G4C replaces deterministic attachment with effective clutches that bind, load, rupture, and
rebind while the cell remains fixed. Actin flow slows with site load, and relative motion loads a
Hookean clutch:
\keylabel
\begin{equation}
 v_a=v_0\max\!\left(0,1-\frac{F_{\mathrm{site}}}{F_{\mathrm{stall}}}\right),
 \qquad
 \dot x=\max(0,v_a-v_{\mathrm{sub}}),
 \qquad f=k_cx.
 \label{eq:motor}
\end{equation}

The Bell slip-bond law gives a continuous force-dependent hazard rather than a deterministic
``too long'' threshold:
\begin{equation}
 k_{\mathrm{off}}(f)=k_{\mathrm{off}}^0\exp\!\left(\frac{f}{F_b}\right),
 \quad
 P_{\mathrm{off}}=1-\exp[-k_{\mathrm{off}}(f)\Delta t],
 \quad
 P_{\mathrm{on}}=1-\exp(-k_{\mathrm{on}}\Delta t).
 \label{eq:bell}
\end{equation}

In the independent baseline,
\begin{equation}
 F_{\mathrm{site}}=\sum_{j\in\mathrm{bound}}k_cx_j.
 \label{eq:independent}
\end{equation}
One rupture removes only that clutch. The whole attachment persists while at least one clutch is
bound.

In the shared-load alternative, if $i$ of $N$ clutches remain bound,
\keylabel
\begin{equation}
 f_i=\frac{F_{\mathrm{site}}}{i},\qquad
 r_i=i k_{\mathrm{off}}^0
 \exp\!\left(\frac{F_{\mathrm{site}}}{iF_b}\right),\qquad
 g_i=(N-i)k_{\mathrm{on}}.
 \label{eq:shared}
\end{equation}
After a rupture, fewer survivors carry the site load, increasing force per survivor and the Bell
off-rate. Rebinding may rescue the site. Complete failure is the first passage to $i=0$.
The implementation also assumes
\begin{equation}
 K_{\mathrm{site}}=Nk_c\quad(i>0),\qquad
 F_{\mathrm{site}}=K_{\mathrm{site}}x_{\mathrm{site}},
 \label{eq:bundle}
\end{equation}
which holds the fully-bound bundle stiffness until the site is completely lost. This is an
explicit coarse-grained closure, not a molecular identity.

The effective defaults are $N=12$, $k_c=2\unit{nN/\mu m}$,
$k_{\mathrm{on}}=0.055\unit{s^{-1}}$, $k_{\mathrm{off}}^0=0.018\unit{s^{-1}}$,
$F_b=1.5\unit{nN}$, $v_0=0.025\unit{\mu m/s}$, and
$F_{\mathrm{stall}}=8\unit{nN/site}$. They describe effective bundles, not twelve counted
integrin molecules.

\section{G4D: reaction-driven rigid-cell motion}

G4D retains the G4C shared-load mechanics and releases only cell translation and rotation.
\keylabel
\begin{equation}
 \gamma_c\dot{\rv}_c=-\sum_m\Fv_m^{\mathrm{cell\to ECM}}.
 \label{eq:translation}
\end{equation}
The cell receives the equal-and-opposite reaction to the traction it applies to collagen. No
$+x$ force, preferred velocity, or hardcoded polarity probability is supplied. Rotation obeys
\begin{equation}
 \gamma_\theta\dot\phi=
 \sum_m(\mathbf a_m-\rv_c)\times[-\Fv_m^{\mathrm{cell\to ECM}}],
 \qquad
 \gamma_\theta=\chi\gamma_cR^2.
 \label{eq:rotation}
\end{equation}
Heterogeneous collagen geometry, different force paths, and stochastic site lifetimes create a
spatial traction imbalance. The fixed and moving panels share the network and random proposals;
only rigid-body degrees of freedom differ. Contacts are periodically refreshed as the cell and
network move; a bound site retains its material attachment until release.

\section{How to read the visualization}

Collagen is shown as coloured curves with small beads connected by axial bonds. Orange marks the
direct-contact graph class, blue one-hop fibres, and purple two-plus-hop fibres. Gold diamonds or
short connectors are permanent elastic crosslinks. Clutch spokes are bound effective clutches;
red marks are clutch rupture events, not fibre or crosslink breaking.

Thin green lines are displacement vectors, not collagen, force links, or crosslinks:
\keylabel
\begin{equation}
 \Delta\rv_i=\rv_i(t)-\rv_i(0),\qquad
 \rv_i^{\mathrm{display}}=\rv_i(0)+M\Delta\rv_i,
 \qquad M\in\{1,10,50\}.
 \label{eq:display}
\end{equation}
Only the green vector length is magnified. True bead and fibre positions remain at $1\times$.
One-hop mean displacement may be only $10^{-3}$--$10^{-2}\unit{\mu m}$ in a
$180\unit{\mu m}$ full-field view, so local zoom and numerical readouts are required.

\section{Parameter provenance and calibration boundaries}

\begin{longtable}{>{\raggedright\arraybackslash}p{0.25\textwidth}p{0.25\textwidth}p{0.40\textwidth}}
\toprule
Parameter & Baseline or range & Interpretation status \\
\midrule
Domain / fibres & $180\unit{\mu m}$ / 99 & Inherited G2-scale fixture; not image-derived density. \\
Cell radius & $10\unit{\mu m}$ & Single-cell fixture; replace with experiment-specific segmentation. \\
Fibre length / diameter & 20--$80\unit{\mu m}$ / $0.30\unit{\mu m}$ & Effective fibre scale, not a molecular fibril. \\
Collagen modulus & $3\unit{MPa}$ & Mechanism-first choice, not a tissue-matched fit. \\
Bending multiplier & 0.10, 0.25, 0.50, 1.0 & Sensitivity separated from axial stiffness. \\
Crosslink probability & 0, 0.20, 0.35, 0.60 & Topology sensitivity; a 2-D crossing is not automatically a biological weld. \\
Crosslink stiffness & 5, 10, 25, $75\unit{nN/\mu m}$ & Effective link-compliance sensitivity. \\
Bead drag & 90, 180, $360\unit{nN\,s/\mu m}$ & Transient sensitivity; requires time-resolved fitting. \\
Total pull & 12, 24, $48\unit{nN}$ & Mechanism range; not fitted cell traction. \\
Contact shell / $\sigma_c$ & $3\unit{\mu m}$ / $1.5\unit{\mu m}$ & Hybrid local-contact starting choice. \\
Cell drag & $600\unit{nN\,s/\mu m}$ & Effective mobility; requires ensemble calibration. \\
\bottomrule
\end{longtable}

\section{Verification gates, limitations, and references}

The required verification gates are: time-step convergence between $0.05$ and
$0.025\unit{s}$; zero traction with no bound clutch; collapse of non-contact transmitted motion
when crosslinks are removed; traceable one-hop force paths; increasing shared-load hazard as
$i$ decreases at fixed site load; complete site loss only at $i=0$; agreement among rupture logs,
traction drops, and animation frames; and paired fixed/moving runs that differ only in cell
degrees of freedom.

G4 is limited to a 2-D projection, one rigid circular cell, linear permanent crosslinks, and
mechanism-first effective parameters. It does not include a nucleus, deformable protrusions,
three-dimensional pores, crosslink turnover, fibre sliding, damage, strain stiffening, SLS
relaxation, proteolysis, proliferation, or biochemical polarity.

\subsection*{Key references}
\begin{enumerate}[leftmargin=1.5em,itemsep=2pt]
\item Lee et al. (2014), \emph{PLOS ONE} 9:e111896. doi:10.1371/journal.pone.0111896.
\item Abhilash et al. (2014), \emph{Biophysical Journal} 107:1829--1840. doi:10.1016/j.bpj.2014.08.029.
\item Wang et al. (2014), \emph{Biophysical Journal} 107:2592--2603. doi:10.1016/j.bpj.2014.09.044.
\item Notbohm et al. (2015), \emph{Journal of the Royal Society Interface}. arXiv:1407.3510.
\item Bell (1978), \emph{Science} 200:618--627. doi:10.1126/science.347575.
\item Bangasser and Odde (2013), \emph{Cellular and Molecular Bioengineering} 6:449--459.
\item Riching et al. (2014), \emph{Biophysical Journal} 107:2546--2558. PMCID: PMC4255204.
\item Erdmann and Schwarz (2004), arXiv:cond-mat/0403552.
\item Gao, Qian, and Chen (2011), \emph{Journal of the Royal Society Interface} 8:1217--1232.
\end{enumerate}

\medskip
\noindent\textbf{Claim boundary.} G4 can separate elastic parameter sensitivity,
crosslink-mediated indirect motion, clutch failure, and reaction-driven rigid-cell motion. It
cannot yet predict whole-tumour invasion.

\end{document}
